\documentclass[11pt]{article}
\usepackage[margin=1in]{geometry}
\Sepackage{amsmath,amssymb,booktabs,graphicx,longtable}
\usepackage[colorlinks=true,linkcolor=blue!50!black,citecolor=blue!50!black]{hyperref}
\usepackage{xcolor}
\newcommand{\new}[1]{\textcolor{blue!60!black}{#1}}
\newcommand{\drop}[1]{\textcolor{red!60!black}{#1}}
\title{\textbf{ProductGPT: revised Model Development and Empirical Results}\\[3pt]
\large Proposed replacement text for Sections~3 and~5 of \texttt{jmp\_v11.tex},
with a change log}
\author{Technical note for co-author review}
\date{30 September 2026}

\begin{document}
\maketitle

\begin{abstract}
\noindent This note proposes replacement text for two sections of
\texttt{jmp\_v11.tex}. The model section is rewritten because the architecture
described there is a subset of what we now estimate: the paper's current text
has no stock variable, and the stock is where the marketing content of the model
lives. The results section is rewritten because its central empirical claim ---
that ProductGPT outperforms recurrent benchmarks --- did not survive a
235-run tuning programme in which every family was tuned to its own frontier.
Section~\ref{sec:log} lists every change and its evidence. Three claims in v11
are \drop{withdrawn}; two reviewer comments are now answerable with data rather
than argument.
\end{abstract}

\section{Change log}\label{sec:log}

\begin{longtable}{p{3.4cm}p{5.2cm}p{5.6cm}}
\toprule
Where & v11 says & Revision, and why \\
\midrule
\endhead
\S3 Model overview & A Transformer over triplet input tokens. & \new{Adds an
explicit \emph{stock} variable and states the model as two memories: a
decision-level memory (the sequence encoder) and a stock-level memory (the
inventory). The stock is where satiation and substitution live, and v11 has no
notation for it.} \\
\addlinespace
\S3 (new) Stock & absent & \new{$S_{it}(j)=\sum_{\tau<t}\kappa(j,c_\tau)\,
g(k_\tau)\,\varphi(t-\tau)$, with the nesting result: $\kappa=I$ and geometric
$\varphi$ give Guadagni--Little loyalty; $\kappa$ = attribute similarity gives
McAlister accumulated-attribute satiation. Attention supplies a \emph{learned}
$\kappa$ and a \emph{learned} $\varphi$.} \\
\addlinespace
\S3 Attention & ``the core innovation \dots\ captures long-range dependencies'' &
\new{Kept, but scoped: the encoder is now one of several interchangeable
choices, and the empirical section shows the choice does not matter for
accuracy. Reviewer Comment~1 (cross-attention vs.\ triplets) is answered with a
measurement rather than a preference.} \\
\addlinespace
\S3 Heterogeneity & absent & \new{Per-customer embedding and a mixture over
output projections, both estimated. Answers the first half of reviewer
Comment~2.} \\
\addlinespace
\S3 Stopping rule & absent & \new{States that NotBuy is an ordinary predicted
class, not an end-of-sequence token, and that it is defined by a fixed 24-hour
interval. Answers the second half of Comment~2.} \\
\addlinespace
\S3 Teacher forcing & ``Without teacher forcing, simulation errors would
accumulate rapidly'' & \new{Correct, and now the starting point of a test rather
than an aside: training on data prefixes while deployment visits model prefixes
is exactly the gap Section~\ref{sec:gen} measures. Answers the third half of
Comment~2.} \\
\addlinespace
\S5 Split & 80/10/10 by customer; campaign 28 = new-LTP holdout, 29--30 =
revived & \new{Replaced by the Lu \& Kannan $2\times2$: customers $\times$
periods, four cells reported separately, validation from late calibration
(campaign 27), \texttt{split\_seed} fixed at 33, holdout opened once.} \\
\addlinespace
\S5 Primary metric & AUC & \new{Held-out negative log-likelihood is the
selection metric; AUC, macro-F1 and hit rate are reported per decision class.
Answers reviewer Comment~3 on metric--goal matching.} \\
\addlinespace
\S5 Model comparison & ``ProductGPT outperforms the RNN benchmarks'' &
\drop{Withdrawn.} \new{Under equal tuning budgets, four families span 0.0194
nats against a seed sd of 0.0075 --- a tie. The best cell is also the smallest
model.} \\
\addlinespace
\S5 Perception map & t-SNE of product embeddings shows category clusters and
figure--weapon pairs & \drop{Withdrawn.} \new{A decomposition shows the learned
identity branch is at chance (0.404 vs.\ 0.408); the apparent structure is the
attribute projection. Confirmed by training with the identity embedding
removed, which costs 0.0001 nats.} \\
\addlinespace
\S5 Triplet vs.\ duplet & AUC 0.9004 vs.\ 0.7611 & \drop{Not reportable as
stated.} \new{Those runs predate two data corrections (a label leak through the
obtained-product block, and a product-index error affecting 89 of 122
products). The ablation must be recomputed on the corrected data.} \\
\addlinespace
\S5 Generative deployment & described as a working decision-support procedure &
\new{Rewritten as an open question with a test attached
(Section~\ref{sec:gen}). Preliminary results do not support the claim as v11
states it.} \\
\bottomrule
\end{longtable}

\section{Proposed replacement: Model Development}

\subsection*{Overview: two memories}

ProductGPT predicts one of $J=9$ decisions at each purchase occasion from three
observed streams: the limited-time products on offer, $\mathbf{x}_{it}$; the
products previously obtained, $\mathbf{o}_{i,t-1}$; and the previous decision,
$y_{i,t-1}$. What the model does with them is best stated as \emph{two
memories}, because they carry different economics and, as
Section~\ref{sec:results} shows, contribute differently.

The \textbf{decision-level memory} summarises the customer's own history of
choices. This is the sequence encoder, and it may be self-attention with a
recency prior, a gated recurrent network, or a hybrid of the two.

The \textbf{stock-level memory} summarises what the customer \emph{holds}. It is
a vector $S_{it}\in\mathbb{R}^{P}$ over the $P$ products, and it is the object
that carries satiation and substitution. The v11 text has no notation for it,
which is why the model there reads as a generic sequence model rather than a
demand model.

\subsection*{The stock, and what it nests}

Let $c_\tau$ index a product acquired at occasion $\tau$, and let $k_\tau$ be
which copy of that product it is (first, second, third). The stock relevant to
product $j$ at occasion $t$ is
\begin{equation}
S_{it}(j)\;=\;\sum_{\tau<t}\ \kappa(j,c_\tau)\ g(k_\tau)\ \varphi(t-\tau),
\label{eq:stock}
\end{equation}
with three primitives:
\begin{itemize}
\item $\varphi(\Delta)=\rho^{\Delta}$, a decay with half-life $h$, so recent
acquisitions matter more;
\item $g(k)$, the weight of a $k$th duplicate copy, normalised to $g(1)=1$;
\item $\kappa(j,p)$, a \textbf{substitution kernel}: how much holding $p$
satiates demand for $j$.
\end{itemize}

Equation~\eqref{eq:stock} is not new to marketing; it is a generalisation of two
classical constructions, and saying so is the cleanest way to explain what the
architecture contributes.

\begin{itemize}
\item With $\kappa=I$ and $\varphi$ geometric, $S_{it}(j)$ is exactly Guadagni
and Little's (1983) exponentially smoothed loyalty variable.
\item With $\kappa$ equal to an attribute-similarity matrix, it is McAlister's
(1982) accumulated-attribute satiation, in which consuming one item of an
attribute class satiates the class.
\end{itemize}

\noindent\textbf{The architecture therefore adds no mechanism; it estimates what
the classics fix in advance.} Attention supplies a learned $\kappa$ --- a
softmax over product embeddings --- and the recency prior supplies a learned
$\varphi$. We verify this nesting by simulation in ordinary brand choice, with
no probabilistic goods anywhere: when a Guadagni--Little process generates the
data the smoothing constant is recovered (0.499, 0.708 and 0.919 against true
0.5, 0.7 and 0.9), and when a McAlister process generates it the attribute
structure is found.

We estimate three forms of $\kappa$: the identity; an attribute
parameterisation with a handful of coefficients; and a free product-level kernel
$\kappa=\mathrm{softmax}(EW_Q(EW_K)^\top/\sqrt{d})$. Section~\ref{sec:results}
reports what each is worth, and a companion note reports which are identified.

\subsection*{Decision-level memory, and the cross-attention question}

Reviewer Comment~1 asks whether cross-attention should pool the three streams
instead of concatenating them into triplets. We implemented both. The
offer-inventory cross-attention --- in which the current offer token attends to
the full history of obtained products --- is the architecture the comment
describes, and it is worth approximately nothing: removing the entire stock
path, cross-attention included, costs 0.0137 nats, inside the seed noise of the
comparison.

We therefore keep the triplet structure, and we say why in the language the
comment invites: the choice is \emph{intentional and empirically grounded},
not a convenience. Cross-attention costs $O(S^2)$ in sequence length --- the
tensor it forms is $(B,H,S,4,10S)$ --- and buys no accuracy on this data.

\subsection*{Product representation}

Each product token is represented as $e_p=\mathrm{id}_p+\gamma\,W\mathbf{f}_p$,
combining a learned identity embedding with a projection of the product's
observed attributes $\mathbf{f}_p$ (rarity, type, element, weapon class and
so on). Section~\ref{sec:results} reports the decomposition of what each branch
contributes, and the answer is consequential enough to change how the paper
discusses embeddings.

\subsection*{Consumer heterogeneity}

Reviewer Comment~2 asks how heterogeneity across players is accounted for. Two
mechanisms are estimated: a per-customer embedding appended to the input, and a
mixture over $H$ output projections with per-customer weights. For customers not
seen in training, the embedding is set to the mean of the trained rows. Both are
evaluated in Section~\ref{sec:results}, and the result is not the one we
expected.

\subsection*{The stopping rule}

Comment~2 also asks how stopping is modelled. Decision 9, ``NotBuy'', is
\emph{not} an end-of-sequence token: it means the customer was shown the
banners and chose not to pull, and they may pull again at the next occasion. It
is an ordinary predicted class, and the majority one. Occasions are defined on a
fixed 24-hour grid --- one NotBuy row is inserted at the end of every interval
in the customer's active window --- so ``NotBuy'' means \emph{no purchase within
one 24-hour interval}. Sequences terminate by padding, which is why no EOS token
is needed. This also means the occasion grid is partly exogenous: the interval
skeleton comes from the calendar, and only the purchase rows are endogenous.

\subsection*{Objective, teacher forcing, and the recursive loop}

The model is trained by minimising cross-entropy on the decision only, with
teacher forcing: after each prediction the \emph{ground-truth} decision, not the
predicted one, is appended to the history. Offers are excluded from the
prediction targets because they are managerial inputs, and outcomes are excluded
because they are drawn by a published lottery.

Comment~2's third question --- how simulation errors are managed in the
recursive generative loop --- is the sharpest of the three, and v11 answers it
only by asserting that teacher forcing prevents the problem during training. It
does, but it also creates one at deployment. Training evaluates the model on
prefixes drawn from the \emph{data}; free-running generation visits prefixes
drawn from the \emph{model}. If the model's conditional distributions were
correctly specified this would not matter --- the chain rule would deliver the
joint law --- so any gap between one-step accuracy and generated-sequence
quality is evidence of misspecification. Section~\ref{sec:gen} measures it
rather than assuming it away.

\section{Proposed replacement: Empirical Results}\label{sec:results}

\subsection*{Design}

We use the Lu and Kannan (2025) $2\times2$: customers are split in half, the
period is split into a calibration block (campaigns 1--27) and a holdout block
(campaigns 28--30), and the four cells are reported separately. The customer
partition is fixed by a dedicated seed so that multi-seed runs share one
partition and their differences reflect training noise. Validation comes from
\emph{late calibration} (campaign 27) rather than from a random slice of time,
so early stopping can see temporal drift; the holdout block is untouched by
model selection and was opened exactly once. The headline cell is out-of-sample
customers $\times$ holdout campaigns.

Selection uses held-out negative log-likelihood. We additionally report AUC,
macro-F1, hit rate and revenue error per decision class, which answers
Comment~3: AUC is reported per binary class and hit rate for the categorical
task, rather than one metric standing in for both.

\subsection*{Tuning protocol}

Because the paper's claim is now about \emph{families} rather than about one
model, the comparison has to be budget-fair. We ran a five-stage programme of
235 runs: a coarse stage over depth and width; a refinement stage sampled from
the leaders; confirmation on fresh seeds; a final stage of five seeds per frozen
configuration; and a diagnostic stage. Every family received an equal budget,
selection used validation only, and a family was adopted only if it beat the
incumbent by more than 0.02 nats on at least four of five paired seeds.

\subsection*{Result 1: no family wins}

\begin{center}\small
\begin{tabular}{llrr}
\toprule
Rank & Family & Holdout NLL & Parameters \\
\midrule
1 & plain GRU & \textbf{0.8856 $\pm$ 0.0067} & 1.25M \\
2 & GRU encoder & 0.8904 $\pm$ 0.0031 & 4.01M \\
3 & transformer + recency prior & 0.8925 $\pm$ 0.0058 & 7.93M \\
4 & hybrid & 0.8958 $\pm$ 0.0076 & 1.47M \\
\bottomrule
\end{tabular}
\end{center}

\noindent The full spread across all eight cells (four families $\times$ two
vocabularies) is 0.0194 nats against a typical seed sd of 0.0075. No pair
satisfies the adoption rule. \textbf{The best cell is also the smallest model}:
1.25M parameters against the transformer's 7.93M, a tie at one sixth of the
size.

This is a stronger statement than v11's, not a weaker one. ``We beat an LSTM''
invites the question of whether the LSTM was tuned. ``Four families, each tuned
to its own frontier under an equal budget, land within 0.02 nats on a holdout
they never saw'' does not.

\subsection*{Result 2: the components are worth less than expected}

\begin{center}\small
\begin{tabular}{lrr}
\toprule
Component removed or added & $\Delta$ validation NLL & Verdict \\
\midrule
stock path removed entirely & $+0.0137$ & inside the margin \\
per-customer embedding added & $+0.0046$ & no gain \\
mixture over output projections added & $+0.0000$ & exact tie \\
both added & $+0.0073$ & no gain \\
product identity embedding removed & $+0.0001$ & no cost \\
product attributes removed & $+0.0064$ & a real cost \\
\bottomrule
\end{tabular}
\end{center}

\noindent Two readings deserve emphasis.

\textbf{The learned identity embedding contributes nothing.} Deleting it costs
0.0001 nats and is slightly better on the holdout; deleting the attributes costs
sixty times as much. Everything the model knows about a product, it knows from
the attribute projection. This is why the perception map in v11 is withdrawn: a
decomposition of the two branches scores the learned identity part at chance
(0.404 against 0.408 for a random baseline), and the visible clustering is the
attribute projection rendered through t-SNE. We report the decomposition instead
of the map.

\textbf{Modelling heterogeneity does not improve prediction.} This is a negative
result on the real panel, and it is reported as such. Its significance is that
the same component is decisive for whether the model's parameters can be
interpreted at all: in simulation, unmodelled taste heterogeneity collapses the
recovered attribute structure and biases the satiation coefficient by a factor
of 2.2, and modelling it removes both problems. A criterion that selects on
predictive fit would have discarded it.

\subsection*{Result 3: the decay is short, and its optimum is interior}

Freezing the half-life and profiling the likelihood over it gives validation NLL
of 0.8795, \textbf{0.8735}, 0.8756, 0.8801 and 0.8817 at half-lives of 0.5, 1,
2, 5 and 10 occasions. The optimum is interior --- 0.5 is worse than 1 and the
curve rises monotonically thereafter --- so \textbf{holdings fade over days, not
over campaigns}. The gap between 1 and 2 is a quarter of a seed sd, so we report
an interval (one to two occasions) rather than a point estimate.

\paragraph{A disclosed limitation.} The runs behind this profile were submitted
with campaign fixed effects enabled, but a defect in the data pipeline meant the
campaign index never reached the model, so the effects were never estimated. We
discovered this on 30 September 2026 and have fixed it. Simulation indicates
that omitting this control \emph{inflates} the estimated half-life, which is the
opposite direction to what we observe, so the qualitative conclusion is not
threatened; nonetheless the profile is being recomputed with the control in
force, and the interval above should be treated as provisional until then.

\subsection*{Result 4: where the families differ, they differ compositionally}

The aggregate tie conceals a systematic pattern. Decomposing the per-class
likelihood, the attention-based encoder is better on \emph{purchase} classes by
0.042 nats, while the recurrent encoder is better on \emph{NotBuy} by 0.079.
The two effects are of comparable size and opposite sign, which is why the
aggregate is flat. The advantage of attention is concentrated in the
Figure-B banner --- the one that is absent in roughly half of campaigns, and
therefore the one where a flexible memory of what was offered when is worth
most.

This is the honest version of the architectural claim: attention does not
predict better overall, but it predicts differently, and where it helps is
interpretable.

\section{Generative validity: an open question}\label{sec:gen}

Section~5 of v11 ends by describing the generative deployment --- the manager
specifies an offer schedule, the model simulates decisions, the published
lottery resolves outcomes --- as a working procedure. That description is
premature, and revising it is the most consequential change proposed here.

Prediction and generation are different claims. Everything above is a statement
about the \emph{next} decision given a true history. A decision-support tool
makes statements about \emph{sequences} under a schedule that has not occurred.
We now test the second directly: condition on each customer's true history
through the calibration period, then roll the model forward over the holdout
campaigns on its own sampled decisions, with the published gacha rules resolving
acquisitions, and compare the resulting sequences to the real ones on statistics
the model was never fitted to --- spend, purchase incidence, run lengths, banner
switching, autocorrelation, and the probability of buying on a banner given how
many copies of its featured product the customer already holds.

The comparison is made against a floor built by splitting the real data against
itself, and against two deliberately weak generators fitted on calibration only:
independent draws from the empirical decision distribution, and a first-order
Markov chain on the nine decisions.

\paragraph{The finding.} No trained configuration matched the real sequences.
Over 300 held-out customers with eight simulated replicates each, all four
configurations --- which are statistically tied on one-step accuracy --- score
between 1.91 and 2.37 in energy distance against a split-half null band of
0.078. A first-order Markov chain scores 0.71, closer to the data than any
trained model, and independent draws score 2.19.

Two conclusions follow, and the second depends on a control.

\textbf{First, the tie survives.} The four configurations differ by 0.0137 nats
in one-step accuracy and by 0.47 in sequence discrepancy against seed standard
deviations of 0.10--0.42. They are indistinguishable on both criteria. Our prior
was that the stock path would separate here, because satiation is a
sequence-level mechanism; it does not.

\textbf{Second, the failure is the feedback loop, not the conditional
distribution.} Teacher forcing --- sampling at each occasion while the input
streams stay real --- is the control that separates a genuine failure of
free-running generation from an error in our rollout:

\begin{center}\small
\begin{tabular}{lrrl}
\toprule
Generator & energy distance & $z$ & verdict \
\midrule
teacher-forced & \textbf{0.307} & 1.6 & \textbf{inside the null band} \
free-running & 1.793 & 9.4 & outside \
independent draws & 1.776 & 9.3 & outside \
\bottomrule
\end{tabular}
\end{center}

\noindent Teacher-forced sequences are statistically indistinguishable from real
ones. The per-step distribution and the scoring are therefore sound, and the
entire discrepancy is created by conditioning the model on its own output.

The mechanism is visible in the generated decision shares:

\begin{center}\small
\begin{tabular}{lrrrrr}
\toprule
& NotBuy & Buy10 Fig-B & Buy1 Weapon & Buy10 Weapon & Buy1 Fig-A \
\midrule
observed & 0.503 & 0.021 & 0.023 & 0.012 & 0.176 \
teacher-forced & 0.410 & 0.048 & 0.036 & 0.020 & 0.169 \
free-running & \textbf{0.183} & 0.113 & 0.134 & \textbf{0.135} & 0.078 \
\bottomrule
\end{tabular}
\end{center}

\noindent Under free running the model abandons the majority class: NotBuy falls
from 0.503 to 0.183, while purchases on the rarer banners inflate five- to
eleven-fold. A sampled purchase updates the inventory, the updated state is one
the model rarely saw during training, and from there it buys more. This is
error accumulation of exactly the form the training objective cannot penalise,
because the likelihood is only ever evaluated on histories the data generated.

\paragraph{What this means for the paper.} The generative deployment described
in v11 is not currently supported, and the managerial section should be
rewritten around what the evidence licenses. Three tiers should be distinguished
explicitly:

\begin{center}\small
\begin{tabular}{p{2.8cm}p{4.4cm}p{6.0cm}}
\toprule
Claim & Status & Licenses \
\midrule
next-decision prediction & established, and a tie across families & forecasting
the next decision \
sequence generation & \textbf{fails, by a controlled test} & nothing, until the
feedback loop is addressed \
structural interpretation & decay and duplicate weights identified; a free
product-level kernel is not & comparative statics on the identified primitives
only \
\bottomrule
\end{tabular}
\end{center}

\noindent This is a more defensible paper than one that asserts the simulation
works. The contribution becomes a demonstration, with a control, that
\emph{predictive accuracy does not license simulation} --- a result that
transfers well beyond this application, and one the current wave of papers using
learned models as consumer simulators has reason to take seriously.

Remedies exist and are the natural next step: scheduled sampling and related
exposure-bias corrections, or training against a sequence-level criterion rather
than one-step likelihood. We report the diagnosis first because it is the part
that generalises.

\end{document}
